\subsection{Commutators and Special Matrix Classes}

\begin{definition} \label{def:operator-and-matrix-commutator} \hlt{Commutator}
\begin{enumerate}[label=(\roman*)]
\setlength{\itemsep}{0pt}
\item For $S,T\in\LL(\mathcal{V})$, their \hlt{commutator} is $[S,T]\defeq ST-TS$.
\item For $A,B\in\operatorname{Mat}(n,n,\mathbb{F})$, their \hlt{matrix commutator} is $[A,B]\defeq AB-BA$.
\end{enumerate}
\end{definition}

\begin{proposition} \label{prop:commutators-have-zero-trace} \hlt{Trace Obstruction for Commutators}\\
Let $A,B,C\in\operatorname{Mat}(n,n,\mathbb{F})$.
\begin{enumerate}[label=(\roman*)]
\setlength{\itemsep}{0pt}
\item $\operatorname{tr}([A,B])=0$.
\item If $C$ is similar to $[A,B]$, then $\operatorname{tr}(C)=0$.
\end{enumerate}
\begin{proof}
\begin{enumerate}[label=(\roman*)]
\setlength{\itemsep}{0pt}
\item $\operatorname{tr}([A,B])=\operatorname{tr}(AB)-\operatorname{tr}(BA)=0$ by Proposition \ref{prop:trace-arithmetic}(i),(ii) and Proposition \ref{prop:trace-cyclicity-two-factors}.
\item If $C$ is similar to $[A,B]$, then $\operatorname{tr}(C)=\operatorname{tr}([A,B])=0$ by Theorem \ref{thm:similarity-preserves-trace-and-determinant}(i) and part (i).
\end{enumerate}
\hlt{Structural Commentary}: Trace cyclicity annihilates a commutator; similarity invariance transfers the obstruction.\\
\hlt{Key Insight}: The two product traces cancel exactly.
\end{proof}
\end{proposition}

\begin{corollary} \label{cor:identity-is-not-operator-commutator} \hlt{The Identity Is Not a Commutator}\\
Let $\mathcal{V}$ be finite-dimensional over $\mathbb{F}$ and $(\dim\mathcal{V})\cdot1_{\mathbb{F}}\neq0$. Then $(\nexists\,S,T\in\LL(\mathcal{V}))\ [S,T]=I_{\mathcal{V}}$.
\begin{proof}
Set $n\defeq\dim\mathcal{V}$, and suppose $[S,T]=I_{\mathcal{V}}$. Choose a basis $\beta$ of $\mathcal{V}$ and set $A\defeq[S]_{\beta\leftarrow\beta}$ and $B\defeq[T]_{\beta\leftarrow\beta}$.\\
Then $[A,B]=[ST-TS]_{\beta\leftarrow\beta}=I_n$ by Propositions \ref{prop:composition-and-matrix-multiplication} and \ref{prop:matrix-representation-respects-linear-combinations}.\\
By Proposition \ref{prop:commutators-have-zero-trace}(i), $0=\operatorname{tr}([A,B])=\operatorname{tr}(I_n)=n\cdot1_{\mathbb{F}}$, contrary to the hypothesis.\\
\hlt{Structural Commentary}: Coordinate matrices transfer the operator equation to the matrix trace obstruction.\\
\hlt{Key Insight}: The identity has trace $(\dim\mathcal{V})\cdot1_{\mathbb{F}}$.
\end{proof}
\end{corollary}

\begin{lemma} \label{lem:scalar-matrix-dependent-orbits} \hlt{Scalar Matrices and One-Step Orbits}\\
Let $n\geq2$ and $A\in\operatorname{Mat}(n,n,\mathbb{F})$. Then $(\forall\,\mathbf{x}\in\mathbb{F}^n)\ (\mathbf{x},A\mathbf{x})$ is linearly dependent iff $A$ is a scalar matrix.
\begin{proof}
\begin{enumerate}[label=\roman*.]
\setlength{\itemsep}{0pt}
\item[$(\Rightarrow)$] Suppose $(\forall\,\mathbf{x}\in\mathbb{F}^n)\ (\mathbf{x},A\mathbf{x})$ is linearly dependent. Since $\mathbf{e}_i\neq\mathbf{0}$, $(\exists\,a_i\in\mathbb{F})\ A\mathbf{e}_i=a_i\mathbf{e}_i$ for $i\in\{1,\ldots,n\}$.\\
For $i\in\{2,\ldots,n\}$, the dependence of $\mathbf{e}_1+\mathbf{e}_i$ and $A(\mathbf{e}_1+\mathbf{e}_i)=a_1\mathbf{e}_1+a_i\mathbf{e}_i$ implies $(\exists\,b_i\in\mathbb{F})\ a_1\mathbf{e}_1+a_i\mathbf{e}_i=b_i(\mathbf{e}_1+\mathbf{e}_i)$.\\
Thus $a_1=b_i=a_i$ by independence of $(\mathbf{e}_1,\mathbf{e}_i)$. Hence $(\forall\,i\in\{1,\ldots,n\})\ A\mathbf{e}_i=a_1\mathbf{e}_i$, so $A=a_1I_n$ by equality on standard basis.
\item[$(\Leftarrow)$] If $A=cI_n$, then $(\forall\,\mathbf{x}\in\mathbb{F}^n)\ A\mathbf{x}=c\mathbf{x}$, so each pair $(\mathbf{x},A\mathbf{x})$ is linearly dependent.
\end{enumerate}
\hlt{Structural Commentary}: Coordinate lines produce individual scalars; diagonal lines force those scalars to coincide.\\
\hlt{Key Insight}: Dependence along $\mathbf{e}_1+\mathbf{e}_i$ equates the coordinate eigenvalues.
\end{proof}
\end{lemma}

\begin{lemma} \label{lem:nonscalar-similar-zero-entry} \hlt{Nonscalar Matrices Have a Zero Diagonal Entry After Similarity}\\
Let $n\geq2$ and $A\in\operatorname{Mat}(n,n,\mathbb{F})$. If $A$ is not scalar, then $A$ is similar over $\mathbb{F}$ to a matrix whose $(1,1)$ entry is $0$.
\begin{proof}
Choose $\mathbf{x}\in\mathbb{F}^n$ such that $(\mathbf{x},A\mathbf{x})$ linearly independent by Lemma \ref{lem:scalar-matrix-dependent-orbits}. Extend this pair to basis $\beta$ of $\mathbb{F}^n$ by Theorem \ref{thm:extending-independent-lists}.\\
Let $T_A:\mathbb{F}^n\to\mathbb{F}^n$ be defined by $T_A(\mathbf{y})\defeq A\mathbf{y}$, and set $C\defeq[T_A]_{\beta\leftarrow\beta}$. The first column of $C$ is $[A\mathbf{x}]_{\beta}=(0,1,0,\ldots,0)^t$.\\
Hence the $(1,1)$ entry of $C$ is zero, and $C$ is similar to $A$ by Theorem \ref{thm:change-of-basis-formula}(i), with the standard basis as the original basis.\\
\hlt{Structural Commentary}: A noninvariant line begins a basis whose second vector is $A\mathbf{x}$.\\
\hlt{Key Insight}: The coordinate vector of $A\mathbf{x}$ begins with zero.
\end{proof}
\end{lemma}

\begin{theorem} \label{thm:similar-to-constant-diagonal} \hlt{Similarity to a Constant-Diagonal Matrix}\\
Let $n\geq1$, $A\in\operatorname{Mat}(n,n,\mathbb{F})$, and $(\forall\,k\in\{1,\ldots,n\})\ k\cdot1_{\mathbb{F}}\neq0$. Then $A$ is similar to a matrix with constant diagonal $\frac{1}{n}\operatorname{tr}(A)$.
\begin{proof}
\begin{enumerate}[label=(\roman*)]
\setlength{\itemsep}{0pt}
\item \hlt{Trace Shift}: Set $\tau\defeq\frac{1}{n}\operatorname{tr}(A)$ and $B\defeq A-\tau I_n$. Then $\operatorname{tr}(B)=\operatorname{tr}(A)-(n\cdot1_{\mathbb{F}})\tau=0$ by Proposition \ref{prop:trace-arithmetic}(i),(ii).\\
It suffices to show that every trace-zero matrix of size at most $n$ is similar to a zero-diagonal matrix.\\
Indeed, adding $\tau I_n$ preserves similarity by Theorem \ref{thm:similarity-preserves-scalar-shifts}(i), with shift $-\tau$.
\item \hlt{Induction}: For $m\in\{1,\ldots,n\}$, let $\mathcal{P}_m$ assert that every trace-zero matrix in $\operatorname{Mat}(m,m,\mathbb{F})$ is similar to a zero-diagonal matrix. Property $\mathcal{P}_1$ holds because a trace-zero $1\times1$ matrix is $[0]$.\\
Fix $m\in\{2,\ldots,n\}$, assume $\mathcal{P}_{m-1}$, and let $C\in\operatorname{Mat}(m,m,\mathbb{F})$ satisfy $\operatorname{tr}(C)=0$.
\item \hlt{Scalar Case}: If $C=cI_m$, then $0=\operatorname{tr}(C)=(m\cdot1_{\mathbb{F}})c$. Thus $c=0$ by $m\cdot1_{\mathbb{F}}\neq0$, so $C$ already has zero diagonal.
\item \hlt{Nonscalar Case}: If $C$ is not scalar, then $C$ is similar to $D\defeq\left[\begin{smallmatrix}0&\mathbf{u}^t\\ \mathbf{v}&C_1\end{smallmatrix}\right]$ by Lemma \ref{lem:nonscalar-similar-zero-entry}.\\
The equality $\operatorname{tr}(C_1)=\operatorname{tr}(D)=\operatorname{tr}(C)=0$ holds by Theorem \ref{thm:similarity-preserves-trace-and-determinant}(i) and the block diagonal entries.\\
Choose $S_1\in\operatorname{Mat}(m-1,m-1,\mathbb{F})^{\times}$ such that $S_1^{-1}C_1S_1$ has zero diagonal by $\mathcal{P}_{m-1}$. Set $R\defeq I_1\oplus S_1$.\\
Then $R^{-1}DR$ has zero diagonal because its lower-right block is $S_1^{-1}C_1S_1$.\\
Thus $C$ is similar to a zero-diagonal matrix by Proposition \ref{prop:similarity-equivalence-relation}, so $\mathcal{P}_m$ holds.\\
Therefore $\mathcal{P}_n$ holds by induction, and the original assertion follows by the trace shift in part (i).
\end{enumerate}
\hlt{Structural Commentary}: A trace shift reduces to zero diagonal; induction isolates one diagonal entry at a time.\\
\hlt{Key Insight}: The remaining block retains trace zero after the first entry vanishes.
\end{proof}
\end{theorem}

\begin{theorem} \label{thm:shoda-trace-zero-commutators} \hlt{Shoda's Theorem in the Distinct-Diagonal Case}\\
Let $n\geq1$, $(\forall\,k\in\{1,\ldots,n\})\ k\cdot1_{\mathbb{F}}\neq0$, and $A\in\operatorname{Mat}(n,n,\mathbb{F})$.
\begin{enumerate}[label=(\roman*)]
\setlength{\itemsep}{0pt}
\item $A$ is a commutator in $\operatorname{Mat}(n,n,\mathbb{F})$ iff $\operatorname{tr}(A)=0$.
\item No nonzero scalar matrix in $\operatorname{Mat}(n,n,\mathbb{F})$ is a commutator.
\end{enumerate}
\begin{proof}
\begin{enumerate}[label=(\roman*)]
\setlength{\itemsep}{0pt}
\item[$(i)$] ($\Rightarrow$) If $A=[U,V]$ for $U,V\in\operatorname{Mat}(n,n,\mathbb{F})$, then $\operatorname{tr}(A)=0$ by Proposition \ref{prop:commutators-have-zero-trace}(i).
\item[$(i)$] ($\Leftarrow$) Suppose $\operatorname{tr}(A)=0$. Choose $S\in\operatorname{Mat}(n,n,\mathbb{F})^{\times}$ such that $B\defeq S^{-1}AS=[b_{ij}]$ has zero diagonal by Theorem \ref{thm:similar-to-constant-diagonal}.\\
For $i\in\{1,\ldots,n\}$, set $d_i\defeq i\cdot1_{\mathbb{F}}$, and let $X\defeq\operatorname{diag}(d_1,\ldots,d_n)$. Define $Y=[y_{ij}]$ by $y_{ii}\defeq0$ and $y_{ij}\defeq(d_i-d_j)^{-1}b_{ij}$ for $i\neq j$.\\
For $i\neq j$, set $k\defeq|i-j|$. Then $1\leq k\leq n-1$ and $d_i-d_j\in\{k\cdot1_{\mathbb{F}},-k\cdot1_{\mathbb{F}}\}\subseteq\mathbb{F}^{\times}$ by the hypothesis.\\
Hence $[X,Y]_{ii}=0=b_{ii}$ and $[X,Y]_{ij}=(d_i-d_j)y_{ij}=b_{ij}$, so $B=[X,Y]$.\\
Therefore $A=[SXS^{-1},SYS^{-1}]$ by direct multiplication, and $A$ is a commutator.
\item[$(ii)$] If $c\in\mathbb{F}\setminus\{0\}$, then $\operatorname{tr}(cI_n)=(n\cdot1_{\mathbb{F}})c\neq0$. Thus $cI_n$ is not a commutator by part (i).
\end{enumerate}
\hlt{Structural Commentary}: Trace supplies necessity; constant-diagonal similarity and distinct scalars construct the converse.\\
\hlt{Key Insight}: Dividing each off-diagonal entry by $d_i-d_j$ realizes the commutator.
\end{proof}
\end{theorem}

\begin{remark} \label{rem:shoda-field-hypothesis} \hlt{Field Hypothesis in the Diagonal Proof}\\
Theorem \ref{thm:shoda-trace-zero-commutators} uses the field hypothesis twice: the scalars $d_1,\ldots,d_n$ must be distinct, and every block size in the induction must remain nonzero. In small characteristic this diagonal construction requires replacement.
\end{remark}

\begin{computationmethod} \label{meth:trace-determinant-invariant-computation} \hlt{Trace--Determinant Invariant Computation}\\
Let $n\geq1$, $A,C\in\operatorname{Mat}(n,n,\mathbb{F})$, and suppose $p_A(z)=\prod\limits_{j=1}^{n}(z-\lambda_j)$ over a field containing $\mathbb{F}$. The following procedure extracts aggregate eigenvalue data and screens a commutator equation:
\begin{enumerate}[label=(\roman*)]
\setlength{\itemsep}{0pt}
\item If $A$ represents an operator, choose a basis in which its representing matrix is easiest to compute.
\item Compute $\operatorname{tr}(A)$ from the diagonal entries and $\det(A)$ by a convenient determinant formula.
\item Record $\sum\limits_{j=1}^{n}\lambda_j=\operatorname{tr}(A)$ and $\prod\limits_{j=1}^{n}\lambda_j=\det(A)$ by Theorem \ref{thm:matrix-characteristic-polynomial-basic-coefficients}(i).
\item For $[X,Y]=C$, first compute $\operatorname{tr}(C)$. If it is nonzero, no solution exists by Proposition \ref{prop:commutators-have-zero-trace}(i).
\end{enumerate}
\end{computationmethod}

\begin{definition} \label{def:idempotent-and-nilpotent-matrices} \hlt{Idempotent and Nilpotent Matrices}\\
Let $A\in\operatorname{Mat}(n,n,\mathbb{F})$.
\begin{enumerate}[label=(\roman*)]
\setlength{\itemsep}{0pt}
\item $A$ is \hlt{idempotent} iff $A^2=A$.
\item $A$ is \hlt{nilpotent} iff $(\exists\,k\in\mathbb{Z}_{>0})\ A^k=0$.
\end{enumerate}
\end{definition}

\begin{proposition} \label{prop:idempotent-matrix-basic-invariants} \hlt{Basic Invariants of Idempotents}\\
Let $P\in\operatorname{Mat}(n,n,\mathbb{F})$ satisfy $P^2=P$, and set $r\defeq\operatorname{rank}(P)$.
\begin{enumerate}[label=(\roman*)]
\setlength{\itemsep}{0pt}
\item $\mathbb{F}^n=\operatorname{range}(P)\oplus\ker(P)$.
\item $(\forall\,\mathbf{y}\in\operatorname{range}(P))\ P\mathbf{y}=\mathbf{y}$ and $(\forall\,\mathbf{z}\in\ker(P))\ P\mathbf{z}=\mathbf{0}$.
\item $P$ is similar to $I_r\oplus0_{n-r}$.
\item $P$ is diagonalizable over $\mathbb{F}$.
\item $\operatorname{spec}(P)\subseteq\{0,1\}$.
\item $\operatorname{tr}(P)=\operatorname{rank}(P)$.
\item If $\mathbb{F}\in\{\mathbb{R},\mathbb{C}\}$, then $P$ is the orthogonal projection onto $\operatorname{range}(P)$ iff $P=P^*$.
\end{enumerate}
\begin{proof}
\begin{enumerate}[label=(\roman*)]
\setlength{\itemsep}{0pt}
\item[$(i)$] For $\mathbf{x}\in\mathbb{F}^n$, write $\mathbf{x}=P\mathbf{x}+(\mathbf{x}-P\mathbf{x})$, where $P\mathbf{x}\in\operatorname{range}(P)$.\\
Since $P(\mathbf{x}-P\mathbf{x})=P\mathbf{x}-P^2\mathbf{x}=\mathbf{0}$, the second term lies in $\ker(P)$.\\
If $\mathbf{x}\in\operatorname{range}(P)\cap\ker(P)$, write $\mathbf{x}=P\mathbf{u}$. Then $\mathbf{x}=P^2\mathbf{u}=P\mathbf{x}=\mathbf{0}$, so the sum is direct.
\item[$(ii)$] If $\mathbf{y}=P\mathbf{u}\in\operatorname{range}(P)$, then $P\mathbf{y}=P^2\mathbf{u}=P\mathbf{u}=\mathbf{y}$. The assertion on $\ker(P)$ follows from its definition.
\item[$(iii)$] Choose a basis $(\mathbf{u}_1,\ldots,\mathbf{u}_r)$ of $\operatorname{range}(P)$ and a basis $(\mathbf{v}_1,\ldots,\mathbf{v}_{n-r})$ of $\ker(P)$ by Corollary \ref{cor:matrix-rank-nullity}(i).\\
Their concatenation is a basis by part (i).\\
By part (ii), the representing matrix in this basis is $I_r\oplus0_{n-r}$. Thus $P$ is similar to this matrix by Theorem \ref{thm:change-of-basis-formula}(i).
\item[$(iv)$] The matrix $I_r\oplus0_{n-r}$ is diagonal, and the similarity in part (iii) is a diagonalization of $P$.
\item[$(v)$] If $P\mathbf{x}=\lambda\mathbf{x}$ for $\mathbf{x}\neq\mathbf{0}$, then $\lambda^2\mathbf{x}=P^2\mathbf{x}=P\mathbf{x}=\lambda\mathbf{x}$. Hence $\lambda(\lambda-1)=0$, so $\lambda\in\{0,1\}$.
\item[$(vi)$] By Theorem \ref{thm:similarity-preserves-trace-and-determinant}(i) and part (iii), $\operatorname{tr}(P)=\operatorname{tr}(I_r\oplus0_{n-r})=r=\operatorname{rank}(P)$.
\item[$(vii)$] ($\Rightarrow$) Regarding $P$ as an operator on the standard inner-product space, $P=P^*$ by Theorem \ref{thm:orthogonal-projection-self-adjoint-idempotent}.
\item[$(vii)$] ($\Leftarrow$) Regarding $P$ as an operator on the standard inner-product space, it is the orthogonal projection onto $\operatorname{range}(P)$ by the idempotence hypothesis, $P=P^*$, and Theorem \ref{thm:orthogonal-projection-self-adjoint-idempotent}.
\end{enumerate}
\hlt{Structural Commentary}: Range--kernel coordinates determine invariants; self-adjointness characterizes orthogonal projection.\\
\hlt{Key Insight}: In an adapted basis, every idempotent is $I_r\oplus0_{n-r}$.
\end{proof}
\end{proposition}
